\documentclass[11pt]{article}
\usepackage[top=.78in,bottom=.76in,left=.85in,right=.85in]{geometry}
\usepackage{amsmath,amssymb,amsthm,booktabs,microtype,xcolor}
\usepackage[colorlinks=true,urlcolor=blue!45!black]{hyperref}
\newtheorem{theorem}{Theorem}
\newcommand{\Z}{\mathbb Z}
\newcommand{\Q}{\mathbb Q}
\begin{document}
\begin{center}
{\LARGE Six-term exclusions for the $-87$ singular modulus}\\[5pt]
{\large A stronger uniform residue certificate}\\[6pt]
{\small 27 November 2026}
\end{center}
\begin{abstract}
For discriminant $-87$, the ProveIt article certifies linear independence
for $3{,}892{,}897$ of the $9{,}657{,}648$ anchored six-exponent residue
multisets modulo $63$. By retaining column-specific 2-adic precision
and testing the resulting lattice modulo $16$, we certify
$8{,}574{,}847$ multisets. Every lift in each certified multiset is covered.
The remaining $1{,}082{,}801$ multisets are not decided here.
\end{abstract}
\begin{theorem}
Let $\alpha$ be a singular modulus of discriminant $-87$. Of the
$9{,}657{,}648$ weakly increasing lists of six integers in $[0,62]$
with smallest entry $0$, at least $8{,}574{,}847$ have the following property:
if six distinct nonnegative exponents have that residue multiset,
their powers of $\alpha$ are linearly independent over $\Q$.
The certificate covers $88.788\%$ of these residue multisets.
\end{theorem}

\paragraph{Setup.}
Use the ring-class polynomial $H=H_{-87}$ of degree six, as in
\href{https://github.com/VladimirReshetnikov/ProveIt/blob/f6b5955e13bcefb2760480e7b823eba55ed2383e/Analysis/GeneralizedRiemannHypothesis/Article/sections/smd.tex}{the merged ProveIt article}
(theorems and questions tagged \texttt{smd}). Let
$A=\Z_{(2)}[X]/(H)$, with $x$ the class of $X$ and basis
$1,x,\ldots,x^5$. Exact reduction of $x^{63}$ modulo $256$ gives
\[
 x^{63}=1+16D,\qquad D\equiv6+9x+3x^2+7x^3+8x^4\pmod {16A}.
\]
For reference, the coefficients of $H$ in ascending order modulo
$16$ are $(1,9,5,8,14,3,1)$. The full input polynomial and
modular reduction are in \texttt{input.py}. An arbitrary six-term
relation can be shifted by its smallest exponent because $\alpha\ne0$,
so it suffices to consider residue multisets containing $0$.

\paragraph{Column precision from Newton divided differences.}
Inside one residue block $a$, write the exponents as $a+63q_i$
with distinct integers $q_i\ge0$. For $0\le j<n_a$, replace the
corresponding powers by the order-$j$ divided difference at the first
$j+1$ of these nodes, multiplied by $j!/16^j$. This is an invertible
triangular column transformation over $\Q$. By the finite binomial
expansion, the new column of order $j$ is
\[
 x^a D^j+
 x^a\sum_{h>j} \frac{j!\,16^{h-j}}{h!}\,
        [q_0,\ldots,q_j]\bigl(T(T-1)\cdots(T-h+1)\bigr)\,D^h .
\]
Divided differences of integer polynomials at integer nodes are
integers, as follows by expressing the divided differences of a monomial
as complete homogeneous symmetric polynomials in the nodes.
Thus a summand has valuation at least $4(h-j)-v_2(h!/j!)$.
For $h=j+1$ this is $4-v_2(j+1)$. For $h-j\ge2$, Legendre's
formula gives a lower bound $3(h-j)-1\ge5$, since the binary digit
sum of $j\le5$ is at most two and that of $h>j$ is positive.
The resulting column is congruent to $x^aD^j$ modulo $2^{s_j}A$,
where
\[
             (s_0,s_1,s_2,s_3,s_4,s_5)=(4,3,4,2,4,3).
\]
In particular, these precisions hold uniformly in all quotient nodes.

\paragraph{A lattice test valid for every lift.}
Multiply column $j$ within each block by $2^{4-s_j}$. The transformed
matrix is now congruent modulo $16$ to the matrix with corresponding
column $2^{4-s_j} x^aD^j$ in the power basis. Over $\Z/16\Z$ use
complete pivoting: at each stage choose an entry of minimum
2-adic valuation in the remaining block, permute it into the pivot
position, and clear entries below it by row operations.
A zero remaining block before the last pivot is inconclusive.
If every pivot is nonzero modulo $16$, the test certifies
that \emph{every lift} of this matrix to $\Z_{(2)}$ is nonsingular:
perform the same operations on a lift. If pivot $i$ has valuation
$r_i<4$, the entries on and to its right in that row are divisible
by $2^{r_i}$ by the pivot rule. Those to the left are multiples of
$16$. Dividing each row by $2^{r_i}$ gives an upper triangular
matrix with unit diagonal modulo $2$. Thus its determinant is
nonzero over $\Q$. The column scalings and divided differences
preserve rational independence, giving the assertion for the
original powers.

\paragraph{Exhaustive finite count.}
Enumerate $0=a_1\le a_2\le\cdots\le a_6<63$ lexicographically.
For any residue of multiplicity $n_a$, append its columns of
orders $0,\ldots,n_a-1$. The checker replays the original
modulo-two test on the same input and verifies that every formerly
certified multiset still passes. A second modulo-sixteen checker
uses odd pivots only, dividing the remaining block by two and
lowering the modulus whenever the block becomes even.
The two stronger tests agree on every multiset.
\begin{center}
\begin{tabular}{rrr}
\toprule
Distinct residues & Certified & Total\\
\midrule
1 & 1 & 1\\
2 & 234 & 310\\
3 & 14,880 & 18,910\\
4 & 313,800 & 378,200\\
5 & 2,379,690 & 2,789,225\\
6 & 5,866,242 & 6,471,002\\
\midrule
All & 8,574,847 & 9,657,648\\
\bottomrule
\end{tabular}
\end{center}
The predecessor's $3{,}892{,}897$ exclusions and this additional
$4{,}681{,}950$ refer to residue multisets, not a proportion of
integer tuples with respect to a natural-density measure.

\paragraph{Boundary.}
This is an unconditional partial answer to \texttt{smd:q:six},
not a proof that $s(-87)=7$. Failure of the test is inconclusive:
it allows a singular matrix among arbitrary congruent matrices, which
need not be realized by any choice of quotient nodes. The next
question is whether higher precision with those nodes constrained,
or supplementary local tests at other primes, can remove the
remaining multisets.
\end{document}
